% !TEX TS-program = lualatex
% !TEX encoding = UTF-8 Unicode
\documentclass[12pt,a4paper]{article}

\IfFileExists{src/libs/body.tex}{% --- body.tex ---
\usepackage{pdfpages}
\usepackage{fontspec}
% Prefer Times New Roman when available; fall back in container environments.
\IfFontExistsTF{Times New Roman}{
  \setmainfont{Times New Roman}
}{
  \setmainfont{Liberation Serif}
}
\usepackage[vietnamese]{babel}

% Page layout and spacing
\usepackage[left=3cm,right=2cm,top=2cm,bottom=2cm]{geometry}
\usepackage{titlesec}
\usepackage{setspace}

% Silence noisy warnings
\usepackage{silence}
\WarningFilter{transparent}{Loading aborted}
\WarningFilter{hyperref}{The anchor of a bookmark and its parent's}

% Core packages
\usepackage{float}
\usepackage{svg}
\usepackage{graphicx}
\graphicspath{{assets/}{../assets/}{../../assets/}}
\svgpath{{assets/}{../assets/}{../../assets/}}
\usepackage{enumitem}
\usepackage{array}
\usepackage{tabularx}
\usepackage{multirow}
\usepackage[table]{xcolor}
\usepackage{ulem}
\usepackage{xurl}

% ====== THÊM 2 GÓI NÀY ĐỂ FIX LỖI ADJUSTBOX VÀ FOREST ======
\usepackage{adjustbox}
\usepackage[edges]{forest}
% ============================================================

% TikZ
\usepackage{tikz}

% ====== BỔ SUNG FIT, BACKGROUNDS, TREES VÀO DANH SÁCH ======
\usetikzlibrary{calc, arrows.meta, shapes.geometric, positioning, fit, backgrounds, trees}

% Table tuning
\renewcommand{\tabularxcolumn}[1]{m{#1}}
\hbadness=10000
\hfuzz=10000pt

% Shared commands
\newcommand{\subsecheading}[1]{%
    \par\vspace{2pt}%
    \hspace{1.5em}\textbf{#1}%
    \par\vspace{2pt}%
}
\newcommand{\introtext}[1]{%
    \par\hspace{3.0em}#1\par%
}
\newcommand{\StandaloneDocumentSetup}[2]{%
  \pagenumbering{arabic}%
  \ApplyStandardFooter{#1}{#2}%
}
\newcommand{\StandaloneSectionSetup}[2]{%
  \ApplyStandardFooter{#1}{#2}%
}

% Math
\usepackage{amsmath}

% Hyperref
\usepackage{hyperref}
\hypersetup{
    colorlinks=true,
    linkcolor=black,
    filecolor=magenta,
    urlcolor=blue,
}}{%
  \IfFileExists{../libs/body.tex}{% --- body.tex ---
\usepackage{pdfpages}
\usepackage{fontspec}
% Prefer Times New Roman when available; fall back in container environments.
\IfFontExistsTF{Times New Roman}{
  \setmainfont{Times New Roman}
}{
  \setmainfont{Liberation Serif}
}
\usepackage[vietnamese]{babel}

% Page layout and spacing
\usepackage[left=3cm,right=2cm,top=2cm,bottom=2cm]{geometry}
\usepackage{titlesec}
\usepackage{setspace}

% Silence noisy warnings
\usepackage{silence}
\WarningFilter{transparent}{Loading aborted}
\WarningFilter{hyperref}{The anchor of a bookmark and its parent's}

% Core packages
\usepackage{float}
\usepackage{svg}
\usepackage{graphicx}
\graphicspath{{assets/}{../assets/}{../../assets/}}
\svgpath{{assets/}{../assets/}{../../assets/}}
\usepackage{enumitem}
\usepackage{array}
\usepackage{tabularx}
\usepackage{multirow}
\usepackage[table]{xcolor}
\usepackage{ulem}
\usepackage{xurl}

% ====== THÊM 2 GÓI NÀY ĐỂ FIX LỖI ADJUSTBOX VÀ FOREST ======
\usepackage{adjustbox}
\usepackage[edges]{forest}
% ============================================================

% TikZ
\usepackage{tikz}

% ====== BỔ SUNG FIT, BACKGROUNDS, TREES VÀO DANH SÁCH ======
\usetikzlibrary{calc, arrows.meta, shapes.geometric, positioning, fit, backgrounds, trees}

% Table tuning
\renewcommand{\tabularxcolumn}[1]{m{#1}}
\hbadness=10000
\hfuzz=10000pt

% Shared commands
\newcommand{\subsecheading}[1]{%
    \par\vspace{2pt}%
    \hspace{1.5em}\textbf{#1}%
    \par\vspace{2pt}%
}
\newcommand{\introtext}[1]{%
    \par\hspace{3.0em}#1\par%
}
\newcommand{\StandaloneDocumentSetup}[2]{%
  \pagenumbering{arabic}%
  \ApplyStandardFooter{#1}{#2}%
}
\newcommand{\StandaloneSectionSetup}[2]{%
  \ApplyStandardFooter{#1}{#2}%
}

% Math
\usepackage{amsmath}

% Hyperref
\usepackage{hyperref}
\hypersetup{
    colorlinks=true,
    linkcolor=black,
    filecolor=magenta,
    urlcolor=blue,
}}{% --- body.tex ---
\usepackage{pdfpages}
\usepackage{fontspec}
% Prefer Times New Roman when available; fall back in container environments.
\IfFontExistsTF{Times New Roman}{
  \setmainfont{Times New Roman}
}{
  \setmainfont{Liberation Serif}
}
\usepackage[vietnamese]{babel}

% Page layout and spacing
\usepackage[left=3cm,right=2cm,top=2cm,bottom=2cm]{geometry}
\usepackage{titlesec}
\usepackage{setspace}

% Silence noisy warnings
\usepackage{silence}
\WarningFilter{transparent}{Loading aborted}
\WarningFilter{hyperref}{The anchor of a bookmark and its parent's}

% Core packages
\usepackage{float}
\usepackage{svg}
\usepackage{graphicx}
\graphicspath{{assets/}{../assets/}{../../assets/}}
\svgpath{{assets/}{../assets/}{../../assets/}}
\usepackage{enumitem}
\usepackage{array}
\usepackage{tabularx}
\usepackage{multirow}
\usepackage[table]{xcolor}
\usepackage{ulem}
\usepackage{xurl}

% ====== THÊM 2 GÓI NÀY ĐỂ FIX LỖI ADJUSTBOX VÀ FOREST ======
\usepackage{adjustbox}
\usepackage[edges]{forest}
% ============================================================

% TikZ
\usepackage{tikz}

% ====== BỔ SUNG FIT, BACKGROUNDS, TREES VÀO DANH SÁCH ======
\usetikzlibrary{calc, arrows.meta, shapes.geometric, positioning, fit, backgrounds, trees}

% Table tuning
\renewcommand{\tabularxcolumn}[1]{m{#1}}
\hbadness=10000
\hfuzz=10000pt

% Shared commands
\newcommand{\subsecheading}[1]{%
    \par\vspace{2pt}%
    \hspace{1.5em}\textbf{#1}%
    \par\vspace{2pt}%
}
\newcommand{\introtext}[1]{%
    \par\hspace{3.0em}#1\par%
}
\newcommand{\StandaloneDocumentSetup}[2]{%
  \pagenumbering{arabic}%
  \ApplyStandardFooter{#1}{#2}%
}
\newcommand{\StandaloneSectionSetup}[2]{%
  \ApplyStandardFooter{#1}{#2}%
}

% Math
\usepackage{amsmath}

% Hyperref
\usepackage{hyperref}
\hypersetup{
    colorlinks=true,
    linkcolor=black,
    filecolor=magenta,
    urlcolor=blue,
}}%
}
\IfFileExists{src/libs/header.tex}{% --- header.tex ---
\usepackage{fancyhdr}
\setlength{\headheight}{15.5pt} % Thêm dòng này để fix cảnh báo chiều cao
%\setlength{\headheight}{14.5pt}
\addtolength{\topmargin}{-2.5pt}

\pagestyle{fancy}
\fancyhf{}
\renewcommand{\headrulewidth}{0pt}
\renewcommand{\footrulewidth}{0pt}

\newcommand{\ApplyStandardHeader}{%
  \fancyhead[L]{%
    \begin{minipage}[t]{0.795\linewidth}%
      UIT \textbar{} SS004.F21 \textbar{} Nhóm 03 \textbar{} Đồ án cuối kỳ%
    \end{minipage}%
    \vrule%
    \begin{minipage}[t]{0.185\linewidth}%
      \centering cTetris%
    \end{minipage}%
  }%
}

\ApplyStandardHeader
}{%
  \IfFileExists{../libs/header.tex}{% --- header.tex ---
\usepackage{fancyhdr}
\setlength{\headheight}{15.5pt} % Thêm dòng này để fix cảnh báo chiều cao
%\setlength{\headheight}{14.5pt}
\addtolength{\topmargin}{-2.5pt}

\pagestyle{fancy}
\fancyhf{}
\renewcommand{\headrulewidth}{0pt}
\renewcommand{\footrulewidth}{0pt}

\newcommand{\ApplyStandardHeader}{%
  \fancyhead[L]{%
    \begin{minipage}[t]{0.795\linewidth}%
      UIT \textbar{} SS004.F21 \textbar{} Nhóm 03 \textbar{} Đồ án cuối kỳ%
    \end{minipage}%
    \vrule%
    \begin{minipage}[t]{0.185\linewidth}%
      \centering cTetris%
    \end{minipage}%
  }%
}

\ApplyStandardHeader
}{% --- header.tex ---
\usepackage{fancyhdr}
\setlength{\headheight}{15.5pt} % Thêm dòng này để fix cảnh báo chiều cao
%\setlength{\headheight}{14.5pt}
\addtolength{\topmargin}{-2.5pt}

\pagestyle{fancy}
\fancyhf{}
\renewcommand{\headrulewidth}{0pt}
\renewcommand{\footrulewidth}{0pt}

\newcommand{\ApplyStandardHeader}{%
  \fancyhead[L]{%
    \begin{minipage}[t]{0.795\linewidth}%
      UIT \textbar{} SS004.F21 \textbar{} Nhóm 03 \textbar{} Đồ án cuối kỳ%
    \end{minipage}%
    \vrule%
    \begin{minipage}[t]{0.185\linewidth}%
      \centering cTetris%
    \end{minipage}%
  }%
}

\ApplyStandardHeader
}%
}
\IfFileExists{src/libs/footer.tex}{% --- footer.tex ---
\usepackage{lastpage}

\newcommand{\ApplyStandardFooter}[2]{%
  \fancyfoot[L]{%
    \begin{minipage}[t]{0.795\linewidth}%
      Document ID: #1%
    \end{minipage}%
    \vrule%
    \begin{minipage}[t]{0.185\linewidth}%
      \centering\thepage/\pageref{#2}%
    \end{minipage}%
  }%
  \fancypagestyle{plain}{%
    \fancyhf{}%
    \renewcommand{\headrulewidth}{0pt}%
    \renewcommand{\footrulewidth}{0pt}%
    \ApplyStandardHeader%
    \fancyfoot[L]{%
      \begin{minipage}[t]{0.795\linewidth}%
        Document ID: #1%
      \end{minipage}%
      \vrule%
      \begin{minipage}[t]{0.185\linewidth}%
        \centering\thepage/\pageref{#2}%
      \end{minipage}%
    }%
  }%
  \pagestyle{fancy}%
}
}{%
  \IfFileExists{../libs/footer.tex}{% --- footer.tex ---
\usepackage{lastpage}

\newcommand{\ApplyStandardFooter}[2]{%
  \fancyfoot[L]{%
    \begin{minipage}[t]{0.795\linewidth}%
      Document ID: #1%
    \end{minipage}%
    \vrule%
    \begin{minipage}[t]{0.185\linewidth}%
      \centering\thepage/\pageref{#2}%
    \end{minipage}%
  }%
  \fancypagestyle{plain}{%
    \fancyhf{}%
    \renewcommand{\headrulewidth}{0pt}%
    \renewcommand{\footrulewidth}{0pt}%
    \ApplyStandardHeader%
    \fancyfoot[L]{%
      \begin{minipage}[t]{0.795\linewidth}%
        Document ID: #1%
      \end{minipage}%
      \vrule%
      \begin{minipage}[t]{0.185\linewidth}%
        \centering\thepage/\pageref{#2}%
      \end{minipage}%
    }%
  }%
  \pagestyle{fancy}%
}
}{% --- footer.tex ---
\usepackage{lastpage}

\newcommand{\ApplyStandardFooter}[2]{%
  \fancyfoot[L]{%
    \begin{minipage}[t]{0.795\linewidth}%
      Document ID: #1%
    \end{minipage}%
    \vrule%
    \begin{minipage}[t]{0.185\linewidth}%
      \centering\thepage/\pageref{#2}%
    \end{minipage}%
  }%
  \fancypagestyle{plain}{%
    \fancyhf{}%
    \renewcommand{\headrulewidth}{0pt}%
    \renewcommand{\footrulewidth}{0pt}%
    \ApplyStandardHeader%
    \fancyfoot[L]{%
      \begin{minipage}[t]{0.795\linewidth}%
        Document ID: #1%
      \end{minipage}%
      \vrule%
      \begin{minipage}[t]{0.185\linewidth}%
        \centering\thepage/\pageref{#2}%
      \end{minipage}%
    }%
  }%
  \pagestyle{fancy}%
}
}%
}
\usepackage{docmute}

\hypersetup{pdftitle={Đặc tả Yêu cầu Phần mềm (SRS) - Phân hệ gameStory}}

\begin{document}
\providecommand{\StandaloneSectionSetup}[2]{\ApplyStandardFooter{#1}{#2}}
\StandaloneSectionSetup{02-gameStoryRequirements.tex}{LastPageGameStoryRequirements}
\onehalfspacing

\section{Đặc tả Yêu cầu Nghiệp vụ (SRS): Phân hệ Cốt truyện (gameStory)}

\subsection{Tổng quan Nghiệp vụ}
Phân hệ \textbf{gameStory} đóng vai trò là màn hình khởi động và giới thiệu của ứng dụng cTetris. Module chịu trách nhiệm thiết lập nhận diện thương hiệu, tải tài nguyên hệ thống, và dẫn dắt người chơi vào không gian trò chơi thông qua cốt truyện tương tác trước khi chuyển sang Menu chính (\texttt{Máy\_Ảo\_Trò\_Chơi}).

\subsection{Mô hình hóa Hệ thống}

\subsubsection{Sơ đồ Ngữ cảnh}
\begin{center}
\begin{adjustbox}{max width=\linewidth}
\begin{tikzpicture}[
    entity/.style={rectangle, draw=black, fill=blue!10, thick, minimum width=2.8cm, minimum height=1.2cm, align=center, font=\small},
    system/.style={circle, draw=orange, fill=orange!10, thick, minimum size=2.8cm, align=center, font=\small},
    arrow/.style={-{Stealth[scale=1.2]}, thick}
]
\node[system] (gamestory) {\textbf{gameStory}};
% Thu hẹp khoảng cách về 3cm để tránh tràn trang A4
\node[entity, left=3cm of gamestory] (player) {Người chơi};
\node[entity, right=3cm of gamestory] (os) {API Cloud/\\OS/ Web};
\node[entity, below=1.5cm of gamestory] (gameconsole) {Phân hệ\\ Máy ảo};
% Tách biệt luồng Người chơi <-> gameStory bằng yshift
\draw[arrow] ($(player.east)+(0,0.3)$) -- node[above, font=\scriptsize, align=center] {Chạy \&\\ Tương tác} ($(gamestory.west)+(0,0.3)$);
\draw[arrow] ($(gamestory.west)-(0,0.3)$) -- node[below, font=\scriptsize] {Hiện UI/Story} ($(player.east)-(0,0.3)$);
% Tách biệt luồng gameStory <-> API Cloud bằng yshift
\draw[arrow] ($(gamestory.east)+(0,0.3)$) -- node[above, font=\scriptsize] {Xin T.Nguyên} ($(os.west)+(0,0.3)$);
\draw[arrow] ($(os.west)-(0,0.3)$) -- node[below, font=\scriptsize] {Cấp Media} ($(gamestory.east)-(0,0.3)$);
% Luồng gameStory -> gameConsole
\draw[arrow] (gamestory) -- node[right, font=\scriptsize] {Báo kết thúc} (gameconsole);
\end{tikzpicture}
\end{adjustbox}
\end{center}

\subsection{Yêu cầu Phiên bản V1: Khởi tạo \& Nhận diện Thương hiệu}
\begin{itemize}[leftmargin=*]
    \item \textbf{Mục tiêu cốt lõi:} Thiết lập nền tảng hoạt động độc lập và nhận diện thương hiệu cTetris thông qua hiệu ứng thị giác 8 giây tối ưu.
    \item \textbf{Đặc tả Giao diện \& Hiệu ứng:}
    \begin{itemize}[noitemsep]
        \item Bố cục dọc tỷ lệ 9:16; Logo cTetris (3$\times$3 ô vuông) Hiện dần chính giữa.
        \item Hiệu ứng văn bản mỏng (Thinning) cho tên trò chơi; Thanh tải màu xanh lá (180$\times$12 px).
    \end{itemize}
    \item \textbf{Luồng tương tác sử dụng:}
    \begin{center}
    \begin{adjustbox}{max width=\linewidth}
    \begin{tikzpicture}[node distance=1cm, every node/.style={font=\scriptsize}, start/.style={circle, draw, fill=green!20, minimum size=0.6cm}, end/.style={circle, draw, fill=red!20, double, minimum size=0.6cm}, process/.style={rectangle, draw, fill=orange!10, minimum width=2cm}, decision/.style={diamond, draw, fill=blue!5, aspect=2}]
        \node (s) [start] {Vào};
        \node (p1) [process, below=0.6cm of s] {Khởi tạo Giờ=0};
        \node (p2) [process, below=0.6cm of p1] {Cập nhật UI};
        \node (d1) [decision, below=0.6cm of p2] {T $\geq$ 8s?};
        \node (e) [end, right=1.5cm of d1] {Ra};
        \draw [->] (s) -- (p1); \draw [->] (p1) -- (p2); \draw [->] (p2) -- (d1);
        \draw [->] (d1) -- node[above] {Có} (e); \draw [->] (d1.west) -- ++(-0.5,0) |- (p2.west);
    \end{tikzpicture}
    \end{adjustbox}
    \end{center}
    \item \textbf{Sơ đồ tình huống sử dụng:}
    \begin{center}
    \begin{adjustbox}{max width=\linewidth}
    \begin{tikzpicture}[node distance=1cm, every node/.style={font=\scriptsize}, actor/.style={circle, draw, fill=gray!10, minimum size=0.5cm}, usecase/.style={ellipse, draw, fill=white, minimum width=2cm}]
        \draw (0,0) rectangle (4,-2.5);
        \node at (2,-0.3) {\textbf{HệThống 1}};
        \node (p) [actor] at (-1.5, -1.25) {N.Chơi};
        \node (u1) [usecase] at (2,-1) {Xem MởĐầu};
        \node (u2) [usecase] at (2,-1.8) {Nhìn CộtTải};
        \draw [->] (p) -- (u1); \draw [->] (p) -- (u2);
    \end{tikzpicture}
    \end{adjustbox}
    \end{center}
    \item \textbf{Yêu cầu Chức năng:} Tự động chạy phần mở đầu; Lấp đầy Thanh Tải theo tuyến tính; Chuyển sang Menu sau 8 giây.
    \item \textbf{Yêu cầu Phi Chức năng:} Tương thích đa nền tảng (SDL3); Hỗ trợ ký tự Bản quyền (©) UTF-8.
    \item \textbf{Yêu cầu Kỹ thuật:} Sử dụng \texttt{nanosvg}; Cơ chế Khởi tạo Trễ Ảnh (Lazy Init Texture) tối ưu tài nguyên.
    \item \textbf{Tiêu chí thành công:} Hệ thống tự động chuyển Menu đúng 8.0 giây; Không vỡ đồ họa SVG.
\end{itemize}

\newpage
\subsection{Yêu cầu Phiên bản V2: Cốt truyện Tương tác \& Đồng bộ Dữ liệu}

\subsubsection{Tổng quan nâng cấp so với V1}

\begin{table}[H]
\centering
\begin{tabularx}{\linewidth}{|l|X|X|}
    \hline
    \textbf{Khía cạnh} & \textbf{V1 (Cơ bản)} & \textbf{V2 (Nâng cấp)} \\
    \hline
    Chế độ vận hành & Hoàn toàn tự động (hẹn giờ 8 giây) & Hai chế độ: mở đầu + hội thoại\\
    \hline
    Cấu trúc hàm 
    & \texttt{Chạy\_Cốt\_Truyện\_Game()} & + \texttt{mãTruyện, mãChương} \\
    \hline
    Tương tác & Không & Bỏ qua / ENTER / TAB / Nhấn chọn \\
    \hline
    Dữ liệu & Không CSDL & SQLite chỉ đọc (\texttt{hội\_thoại\_chia\_sẻ}, \texttt{lựa\_chọn\_chia\_sẻ}) \\
    \hline
    Mạng lưới & Không & Đồng bộ Gist qua hàm băm $\rightarrow$ CSDL SQLite \\
    \hline
    Thanh tải & Điền theo thời gian & Kết hợp: \texttt{max(Tiến\_độ\_Đồng\_bộ, Tiến\_độ\_Thời\_gian)} \\
    \hline
    Âm thanh & Không & Nhạc nền phụ (lưu URL, 
    phát ở V3) \\
    \hline
    Luồng sau nhập & Tự động $\rightarrow$ Máy Ảo & Kiểm tra điều kiện $\rightarrow$ chạy / hỏi mở khóa \\
    \hline
\end{tabularx}
\end{table}

\subsubsection{Sơ đồ Trạng thái V2 (Cỗ máy Trạng thái Hai bước)}

\noindent \texttt{mãTruyện = 0} (gọi từ \texttt{hàm\_chính}) kích hoạt Bước 1+2.
\texttt{mãTruyện > 0} (từ Máy Ảo) nhảy thẳng vào chế độ Hội Thoại.
\begin{center}
\begin{adjustbox}{max width=\linewidth}
\begin{tikzpicture}[
    node distance=2.0cm,
    state/.style={circle, draw=blue!60, fill=blue!5, thick,
                  minimum size=1.8cm, align=center, font=\scriptsize},
    sync/.style={rectangle, draw=orange!70, fill=orange!5, thick,
                 rounded corners, minimum width=2.2cm,
                 minimum height=0.6cm, align=center, font=\scriptsize},
    arrow/.style={-{Stealth[scale=1.1]}, thick},
    darrow/.style={-{Stealth[scale=1.1]}, thick, dashed, gray}
]
    % Sync states (top row)
    \node[sync] (sidle)   at (0, 3.5)    {CHỜ\_ĐBỘ};
    \node[sync] (sfetch)  at (2.8, 3.5)  {ĐANG\_TẢI\_VỀ};
    \node[sync] (supd)    at (5.6, 3.5)  {ĐANG\_LƯU\_CSDL};
    \node[sync] (sdone)   at (8.4, 3.5)  {XONG\_ĐBỘ};
    \node[sync, draw=red!60, fill=red!5]
               (soff)    at (5.6, 2.4)  {NGOẠI\_TUYẾN};
    \draw[arrow] (sidle)  -- node[above,font=\tiny]{chạy} (sfetch);
    \draw[arrow] (sfetch) -- node[above,font=\tiny]{m.fest tốt} (supd);
    \draw[arrow] (supd)   -- node[above,font=\tiny]{hoàn tất} (sdone);
    \draw[arrow] (sfetch.south) -- node[right,font=\tiny]{lỗi/hết giờ} (soff.north west);
    \draw[arrow] (supd.south)   -- node[right,font=\tiny]{lỗi tải} (soff.north);
    % Phase 2 states (middle row)
    \node[state] (logo)  at (2.0, 0)   {BƯỚC\\LOGO};
    \node[state] (diag)  at (5.6, 0)   {MÁY\_CHẠY\\THOẠI};
    \node[state, draw=green!60, fill=green!5]
                (done)  at (8.8, 0)   {BƯỚC\\XONG};
    % sync -> logo (both done and offline)
    \draw[arrow] (sdone.south) -- ++(0,-0.8) -| node[below,font=\tiny,pos=0.25]{đ.bộ xong} (logo.north);
    \draw[arrow] (soff.south)  -- ++(0,-0.3) -|
        node[below,font=\tiny,pos=0.1]{ko mạng} (logo.north west);
    % logo -> dialogue (post-sync condition check)
    \draw[arrow] (logo.east) -- node[above,font=\tiny,align=center]{8s/Bỏ\\$\rightarrow$ KT\_SauĐB()} (diag.west);
    \draw[arrow] (diag.east) -- node[above,font=\tiny]{HẾT / Bỏ} (done.west);
    \draw[darrow] (logo.east) to[bend right=20]
        node[below,font=\tiny]{kếTiếp=0} (done.south);
    % Dialogue direct entry (storyId > 0)
    \node[font=\scriptsize, draw=blue!30, fill=white, rounded corners,
          inner sep=3pt] (ext) at (5.6, -2.0)
        {Vào thẳng: \texttt{ID\_Truyện>0}};
    \draw[arrow, blue!60] (ext.north) -- (diag.south);
    % Phase labels
    \node[draw=orange!40, dashed, rounded corners, fill=orange!3,
          fit=(sidle)(sfetch)(supd)(sdone)(soff),
          label=left:{\small\textbf{Bước 1}}] {};
    \node[draw=blue!40, dashed, rounded corners, fill=blue!3,
          fit=(logo)(diag)(done),
          label=left:{\small\textbf{Bước 2}}] {};
\end{tikzpicture}
\end{adjustbox}
\end{center}

\subsubsection{Kiến trúc Đồng bộ Dữ liệu V2 (Kho Gist \texorpdfstring{$\rightarrow$}{->} SQLite)}

\noindent V2 dùng Kho Gist của GitHub làm \emph{kênh phân phối nội dung} (chênh lệch mã băm, không thay thế toàn bộ).
V3 sẽ chuyển sang dùng Cloudflare D1 / Durable Objects.

\begin{center}
\begin{adjustbox}{max width=\linewidth}
\begin{tikzpicture}[
    every node/.style={font=\scriptsize},
    box/.style={rectangle, draw=black, thick, fill=white,
                inner sep=5pt, align=left, minimum width=3.0cm},
    arrow/.style={-{Stealth[scale=1.1]}, thick},
    darrow/.style={-{Stealth[scale=1.1]}, thick, dashed}
]
    \node[box, fill=yellow!10] (gist) at (0, 0) {
        \textbf{Kho Gist}\\
        \rule{2.6cm}{0.4pt}\\
        \texttt{bản\_kê.json}\\
        \quad \texttt{mới: "v2"}\\
        \quad \texttt{v2.c001: \{sha, gistUrl\}}\\
        \quad \texttt{v2.c002: \{sha, gistUrl\}}\\
        \texttt{chapter\_c001.json}\\
        \texttt{chapter\_c002.json}
    };
    \node[box, fill=orange!10, right=2.2cm of gist] (sync) {
        \textbf{Đồng\_bộ\_Gist()}\\
        \rule{2.6cm}{0.4pt}\\
        1. TẢI bản\_kê\\
        2. So mã SHA vs\\
        \quad \texttt{siêu\_dữliệu}\\
        3. TẢI file\\
        \quad chương MỚI\\
        4. \texttt{ÁpDụng\_Chương()}\\
        5. \texttt{Lưu\_Mã\_SHA()}\\
        6. Đ.BỘ IDBFS
    };
    \node[box, fill=blue!10, right=2.2cm of sync] (db) {
        \textbf{CSDL (mặc\_định.db)}\\
        \rule{2.6cm}{0.4pt}\\
        \texttt{dữliệu\_k}\\
        \texttt{hộithoại\_k}\\
        \texttt{lựachọn\_k}\\
        \texttt{siêudliệu}\\
        \quad (mã\_chương, sha,\\
        \quad  url\_phươngtiện)
    };
    \draw[arrow] (gist.east) -- node[above]{TảiĐ.BộHTTP} (sync.west);
    \draw[arrow] (sync.east) -- node[above,align=center]{THÊM HAY\\GHI ĐÈ} (db.west);
    \draw[darrow, gray] (db.south) -- ++(0,-0.5)
        node[below, draw=gray!40, fill=gray!5, rounded corners, inner sep=2pt]
        {\scriptsize V3: chuyển qua Cloudflare D1 / Durable Object};
    % SHA check
    \draw[darrow, blue!50] (db.north) to[bend right=35]
        node[above, font=\tiny]{tìm SHA nội bộ} (sync.north east);
\end{tikzpicture}
\end{adjustbox}
\end{center}

\subsubsection{Kiến trúc Bộ Máy Hội Thoại (Phần D)}
\textit{Kiến trúc nội tại của Bộ Máy Hội Thoại (sơ đồ thành phần, luồng dữ liệu,
sơ đồ trình tự giữa \texttt{Thao\_Tác\_Cơ\_Sở\_Dữ\_Liệu\_Truyện*()}, \texttt{Vòng\_Lặp\_Hội\_Thoại()},
\texttt{Vẽ\_Giao\_Diện\_Hội\_Thoại()}) được đặc tả đầy đủ trong
\textbf{SDS §05-gameStoryEngineering, mục 2.3--2.4}.}

Yêu cầu chức năng của Phần D (điều hướng ngõ vào, phân nhánh lựa chọn, điều kiện
kết thúc) được phát biểu trong §\ref{subsec:v2-func} dưới đây.

\subsubsection{Bố cục màn hình Bộ Máy Hội Thoại (270×480)}\label{subsec:v2-func}
\textit{Bố cục phân giải pixel (Khu\_Vực\_Ảnh, hộp chữ, các khung lựa chọn, tọa độ
chính xác NÚT\_BỎ\_QUA) được vẽ trong \textbf{SDS §05-gameStoryEngineering, mục 3.1 (Bố cục màn hình V1)}.
SRS chỉ phát biểu ràng buộc giao diện:}
\begin{itemize}[nosep]
  \item Ảnh bối cảnh chiếm $260\times196\,\text{px}$ vùng trên (y=27..223).
  \item Hộp chữ $260\times244\,\text{px}$ vùng dưới, nền bán trong suốt (\texttt{độ\_mờ=235}).
  \item Nút Bỏ Qua luôn hiển thị ở góc trên phải; ẩn hoàn toàn trong phần đồng bộ.
\end{itemize}

\subsubsection{Luồng Kiểm tra Điều kiện Sau Đồng bộ (\texttt{Kiểm\_Tra\_Điều\_Kiện\_Sau\_Đồng\_Bộ})}

\begin{center}
\begin{adjustbox}{max width=\linewidth}
\begin{tikzpicture}[
    every node/.style={font=\scriptsize},
    proc/.style={rectangle, draw=black, fill=orange!10, thick,
                 minimum width=4.2cm, minimum height=0.55cm, align=center},
    dec/.style={diamond, draw=black, fill=green!10, thick,
                aspect=2.5, align=center},
    term/.style={rectangle, rounded corners, draw=black, fill=red!10, thick,
                 minimum width=3.0cm, minimum height=0.55cm, align=center},
    arrow/.style={-{Stealth[scale=1.1]}, thick}
]
    \node[proc]  (q) {Tìm: \texttt{DS\_TruyệnCủaUser CÓ ĐangChọn=1}};
    \node[dec,   below=0.5cm of q]  (d1) {đã chọn\\truyện?};
    \node[proc,  below=0.5cm of d1] (fs) {dùng truyện số 1 ở \texttt{dữliệu\_k ch=1}};
    \node[dec,   below=0.5cm of fs] (d2) {thoả đ.kiện\\truyện kế?};
    \node[term,  right=1.8cm of d2] (unlock) {trả $-$mãKế\\(hỏi CÓ/KHÔNG)};
    \node[term,  below=0.5cm of d2] (replay) {trả +mãHiệntại\\(chơi lại)};
    \node[term,  right=1.8cm of d1] (skip)   {trả 0\\($\rightarrow$ M.H.Chính)};
    \draw[arrow] (q) -- (d1);
    \draw[arrow] (d1) -- node[left]{ko} (fs);
    \draw[arrow] (d1) -- node[above,align=center]{ko có\\hàng nào} (skip);
    \draw[arrow] (fs) -- (d2);
    \draw[arrow] (d2) -- node[above]{có} (unlock);
    \draw[arrow] (d2) -- node[left]{ko} (replay);
\end{tikzpicture}
\end{adjustbox}
\end{center}

\subsubsection{Sơ đồ Tình huống Sử dụng V2}

\begin{center}
\begin{adjustbox}{max width=\linewidth}
\begin{tikzpicture}[
    every node/.style={font=\scriptsize},
    actor/.style={circle, draw, fill=gray!10, minimum size=0.6cm},
    usecase/.style={ellipse, draw, fill=white, minimum width=2.8cm, align=center}
]
    % Đã mở rộng khung viền một chút sang phải để chứa đường cong
    \draw (0,0) rectangle (7.5,-5.2);
    \node at (3.75,-0.3) {\textbf{TruyệnGame 2}};
    \node (p)   [actor] at (-1.6, -2.5)  {N.Chơi};
    \node (net) [actor, draw=blue!60] at (9.1, -2.5) {Kho / Mạng};
    \node (u1) [usecase] at (3.75,-0.9) {Xem MởĐầu + Bỏ};
    \node (u2) [usecase] at (3.75,-1.8) {Đọc ĐốiThoại};
    \node (u3) [usecase] at (3.75,-2.7) {Chọn Đường (A/B)};
    \node (u4) [usecase] at (3.75,-3.6) {Xem Hỏi MởKhóa};
    \node (u5) [usecase, draw=blue!60] at (3.75,-4.6) {ĐồngBộ BảnKê SHA};
    \draw [->] (p) -- (u1); 
    \draw [->] (p) -- (u2);
    \draw [->] (p) -- (u3); 
    \draw [->] (p) -- (u4);
    \draw [->, blue!60] (net) -- (u5);
    % Hiệu chỉnh: Sử dụng vòng cung (bend right) nối từ rìa phải để không đâm xuyên các khối
    \draw [->, dashed] (u5.east) to[bend right=50] (u2.east);
\end{tikzpicture}
\end{adjustbox}
\end{center}

\subsubsection{Cây tính năng V2 (Cây phân bổ nâng cấp)}

\begin{center}
\begin{adjustbox}{max width=\linewidth}
\begin{tikzpicture}[
    grow=right,
    level 1/.style={sibling distance=3.8cm, level distance=4.5cm},
    level 2/.style={sibling distance=1.3cm, level distance=4.0cm},
    edge from parent fork right,
    every node/.style={rectangle, draw=black, rounded corners,
                       fill=gray!10, thick, align=center, font=\scriptsize, text width=2.2cm},
    root/.style={fill=orange!20, font=\small\bfseries, text width=2.0cm}
]
\node[root] {ĐổiMớiV2}
    child { node[fill=blue!10] {Tầng Đ.Bộ}
        child { node {So mãSHA\\Gist} }
        child { node {ápDụngFile\\Chương()} }
        child { node {theoDõi\_Siêu\\DữLiệu} }
    }
    child { node[fill=yellow!15] {MáyChạy\\HộiThoại}
        child { node {CỗMáyTrạngThái\\NútHộiThoại} }
        child { node {PhânNhánh\\LựaChọn} }
        child { node {NhạcNềnTạm\\TừngNút} }
    }
    child { node[fill=green!10] {XửLý\\SauĐ.Bộ}
        child { node {KiểmTra\\ĐiềuKiện} }
        child { node {HỏiĐáp\\MởKhóa} }
        child { node {LựaChọn\\CốtTruyện} }
    }
    child { node[fill=red!10] {UI V2}
        child { node {Nút BỏQua\\(khi logo)} }
        child { node {Cột tiếnđộ\\đ.bộ} }
    };
\end{tikzpicture}
\end{adjustbox}
\end{center}

\subsubsection{Yêu cầu Chức năng V2}
\begin{itemize}[leftmargin=*]
    \item \textbf{Hai ngõ vào độc lập:} \texttt{mãTruyện=0} $\rightarrow$ mở đầu + đồng bộ;
    \texttt{mãTruyện>0} $\rightarrow$ hội thoại trực tiếp (từ bảng Chọn Truyện trong Máy Ảo). Tương thích ngược với hàm gọi \texttt{hàm\_chính} từ V1.
    \item \textbf{Đồng bộ Gist theo hàm băm SHA:} so sánh \texttt{siêu\_dữ\_liệu.sha} so với bản kê khai; chỉ tải và áp dụng các chương có đoạn băm SHA mới.
    Mất mạng $\rightarrow$ \texttt{TRẠNG\_THÁI\_NGOẠI\_TUYẾN}, tiếp tục dùng cơ sở dữ liệu nội bộ.
    \item \textbf{Cập\_Nhật\_Dữ\_Liệu\_Chương():} XÓA các truyện không còn trong tập tin JSON $\rightarrow$ THÊM HOẶC CẬP NHẬT \texttt{dữ\_liệu\_được\_chia\_sẻ} + \allowbreak\texttt{hội\_thoại\_được\_chia\_sẻ} + \allowbreak\texttt{lựa\_chọn\_được\_chia\_sẻ} trong một giao dịch đơn.
    Không thay thế hoàn toàn để bảo toàn quá trình trải nghiệm của người dùng.
    \item \textbf{Bộ máy hội thoại:} tải hàng loạt các nút $\rightarrow$ cỗ máy trạng thái \texttt{Vòng\_Lặp\_Xử\_Lý\_Hội\_Thoại()};
    các lựa chọn được tải theo yêu cầu tại từng nút. Kết thúc: \texttt{mãNútTiếpTheo=0} $\rightarrow$ thoát.
    \item \textbf{Nút Bỏ qua:} CHỈ hiển thị trong BƯỚC\_LOGO (đồng bộ đã hoàn thành);
    ẩn trong khoảng thời gian đồng bộ để tránh người dùng ngắt tiến trình ghi CSDL.
    \item \textbf{Thanh tải tiến trình V2:} \texttt{độ\_dài\_màu\_nền = max(Tiến\_độ\_Đồng\_bộ, Tiến\_độ\_Thời\_gian)};
    dòng trạng thái hiển thị chương đang đồng bộ / "Ngoại Tuyến" / "Đã Cập Nhật Mới Nhất".
    \item \textbf{Kiểm tra điều kiện sau khi đồng bộ:} trả về 0=Máy Ảo, +N=chơi lại, $-$N=hỏi mở khóa truyện mới.
\end{itemize}

\subsubsection{Yêu cầu Phi Chức năng V2}
\begin{itemize}[noitemsep, leftmargin=*]
    \item \textbf{Giao thức HTTP:} môi trường gốc = libcurl chặn luồng (chấp nhận tải khi vừa mở);
    môi trường Trình duyệt Web = \texttt{tải\_qua\_emscripten} ở chế độ đồng bộ qua luồng bất đồng bộ (\texttt{-sASYNCIFY=1}).
    \item \textbf{SQLite:} Phân hệ cốt truyện chỉ \emph{đọc} bảng \texttt{hội\_thoại/lựa\_chọn};
    việc \emph{ghi} chỉ áp dụng cho \texttt{siêu\_dữ\_liệu} (lưu dấu vết mã băm SHA) và bảng \texttt{truyện\_người\_dùng} (việc lựa truyện). Không can thiệp ghi bảng \texttt{kỷ\_lục\_người\_dùng}.
    \item \textbf{Bảo lưu Trình Duyệt:} \texttt{EM\_ASM(\{Module.FS.syncfs(false,\ldots)\})} ngay sau mỗi khi áp dụng dữ liệu chương và sau khi xác nhận truy cập cốt truyện.
    \item \textbf{Cơ chế Giảm Cấp Cục Bộ:} \texttt{LIÊN\_KẾT\_GIST\_BẢN\_KÊ} rỗng hoặc quá trình tải gặp lỗi $\rightarrow$ \texttt{TRẠNG\_THÁI\_NGOẠI\_TUYẾN} $\rightarrow$ tiếp tục chạy hoạt động bình thường với dữ liệu SQLite đã lưu trên máy tính.
\end{itemize}

\newpage
\subsubsection{Tiêu chí Thành công V2}
\begin{table}[H]
\centering
\begin{tabularx}{\linewidth}{|l|X|}
    \hline
    \textbf{Kịch bản kiểm thử} & \textbf{Kết quả mong đợi} \\
    \hline
    SHA khớp mọi chương & Thanh tải đầy 100\% ngay;
    ghi chú "mọi chương cập nhật"; không yêu cầu HTTP cho tệp JSON của chương \\
    \hline
    SHA đổi ở 1 chương & Dòng trạng thái "Đang đồng bộ c001...";
    dữ liệu cập nhật; mã băm siêu dữ liệu cập nhật; đẩy vào IDBFS \\
    \hline
    Mất mạng (Lỗi MANIFEST\_GIST\_URL) & "Ngoại tuyến - dùng dữ liệu cục bộ; mở đầu bình thường; hội thoại dùng DB đã lưu" \\
    \hline
    mãTruyện>0 từ Máy ảo & Bỏ qua mở đầu/đồng bộ;
    tải thẳng hội thoại; Bỏ qua $\rightarrow$ về Máy ảo \\
    \hline
    Phân nhánh hội thoại (cóLựaChọn=đúng) & TAB chuyển lựa chọn;
    ENTER/click để chọn; mãNútKế tải đúng nhánh \\
    \hline
    mở khóa truyện kế (Kiểm\_Tra\_Sau\_Đồng\_Bộ < 0) & Hỏi CÓ/KHÔNG hiện;
    CÓ $\rightarrow$ Cập nhật CSDL đangChọn + xóa IDBFS + chạy truyện;
    KHÔNG $\rightarrow$ chơi lại hiện tại \\
    \hline
    WASM: đồng bộ chương $\rightarrow$ đóng tab $\rightarrow$ mở lại & Siêu dữ liệu SHA còn;
    bỏ qua đồng bộ chương đã tải \\
    \hline
\end{tabularx}
\end{table}

\subsection{Yêu cầu Phiên bản V3: Tích hợp API \& Lưu trữ Đám mây}
\begin{itemize}[leftmargin=*]
    \item \textbf{Mục tiêu cốt lõi:} Tối ưu dung lượng bằng cách tải tài nguyên Đa phương tiện động từ Đám mây dựa trên tốc độ mạng thực tế.
    \item \textbf{Đặc tả Giao diện \& Hiệu ứng:} Trạng thái Tải xuống (KB/s, \%); Hiệu ứng Logo lặp (Pulse) báo hiệu đang tải.
    \item \textbf{Luồng tương tác sử dụng:}
    \begin{center}
    \begin{tikzpicture}[node distance=1cm, every node/.style={font=\scriptsize}, start/.style={circle, draw, fill=green!20}, end/.style={circle, draw, fill=red!20, double}, process/.style={rectangle, draw, fill=orange!10}, decision/.style={diamond, draw, fill=blue!5, aspect=2}]
        \node (s) [start] {Mở};
        \node (d1) [decision, below=0.6cm of s] {CóMạng?};
        \node (p1) [process, below=0.6cm of d1] {Tải T.Liệu};
        \node (p2) [process, below=0.6cm of p1] {Chạy Truyện(V2)};
        \node (e) [end, right=1.5cm of d1] {Đóng(V1)};
        \draw [->] (s) -- (d1);
        \draw [->] (d1) -- node[left] {Có} (p1); \draw [->] (d1) -- node[above] {Ko} (e);
        \draw [->] (p1) -- (p2); \draw [->] (p2.east) -| (e.south);
    \end{tikzpicture}
    \end{center}
    \item \textbf{Sơ đồ tình huống sử dụng:}
    \begin{center}
    \begin{tikzpicture}[node distance=1cm, every node/.style={font=\scriptsize}, actor/.style={circle, draw, fill=gray!10}, usecase/.style={ellipse, draw, fill=white, minimum width=2.5cm}]
        \draw (0,0) rectangle (5,-3);
        \node at (2.5,-0.3) {\textbf{HệThống 3}};
        \node (sys) [actor] at (6.5, -1.5) {Đám mây};
        \node (u1) [usecase] at (2.5,-1) {Kiểm tra mạng};
        \node (u2) [usecase] at (2.5,-2) {Tải Media};
        \draw [->] (sys) -- (u2);
    \end{tikzpicture}
    \end{center}
    \item \textbf{Yêu cầu Chức năng:} Đồng bộ Thanh tải với tiến độ tải về;
    Tự động Bỏ qua nếu mất mạng; Xóa tài nguyên dữ liệu lớn khỏi file xuất xưởng.
    \item \textbf{Yêu cầu Phi Chức năng:} Cơ chế Hết thời gian chờ (Timeout) sau 15 giây; Bảo mật tài nguyên; Chống treo luồng chính.
    \item \textbf{Yêu cầu Kỹ thuật:} Hàm \texttt{tải\_qua\_emscripten} (Trình duyệt) hoặc Lập trình Mạng (Bản cài gốc); Quản lý Bộ nhớ đệm tài nguyên động.
    \item \textbf{Tiêu chí thành công:} Giảm 50\% dung lượng file .exe; Cốt truyện chỉ chạy sau khi tải xong 100\%.
\end{itemize}

\label{LastPageGameStoryRequirements}
\end{document}